\documentclass[11pt]{article}
\usepackage[a4paper,margin=28mm]{geometry}
\usepackage{amsmath,amssymb,amsthm,mathtools,booktabs}
\usepackage[T1]{fontenc}
\usepackage{lmodern,microtype}
\usepackage[hidelinks]{hyperref}
\newtheorem{theorem}{Theorem}[section]
\newtheorem{lemma}[theorem]{Lemma}
\newtheorem{proposition}[theorem]{Proposition}
\newtheorem{corollary}[theorem]{Corollary}
\theoremstyle{remark}\newtheorem{remark}[theorem]{Remark}
\newcommand{\R}{\mathbb R}
\newcommand{\C}{\mathbb C}
\newcommand{\Z}{\mathbb Z}
\newcommand{\spec}{\operatorname{spec}}
\newcommand{\diag}{\operatorname{diag}}
\newcommand{\tr}{\operatorname{tr}}
\newcommand{\norm}[1]{\left\lVert#1\right\rVert}
\newcommand{\rstar}{r_*}
\newcommand{\tw}{\rho_{\mathrm{tw}}}
\title{Holonomy and spectral bounds for signed cycle squares}
\author{}
\date{Research manuscript increment\quad 5 October 2026}
\begin{document}
\maketitle
\begin{abstract}
We study the adjacency spectral radius of real signings of the cycle square
$C_n(1,2)$. For the triangle-sign word $(+,+,-,+,-,-,+,-)$ repeated on
$8m$ vertices, we give a self-contained proof of the exact radii in the two
Hamilton-cycle holonomy sectors. Negative holonomy has radius
$\sqrt{4+\sqrt{8+2\cos(\pi/m)+\sqrt{26-6\cos(\pi/m)}}}$,
strictly below the positive-holonomy value
$\sqrt{4+\sqrt{10+2\sqrt5}}$ for every finite $m$. This disproves the
unrestricted-minimum conjecture in Suvagiya's September 2026 revision.
Separately, an explicit signing on every $n\equiv2\pmod8$, $n\ge50$,
satisfies $\rho^2<198/25$. Its proof combines a rational finite entrance
certificate with an analytic Riccati contraction and a uniform Schur-tail
bound. Neither result determines the unrestricted minimum over all signings.
\end{abstract}

\section{Introduction}

For $n\ge8$, the cycle square $C_n(1,2)$ has vertex set $\Z/n\Z$ and
edges $\{i,i+1\}$ and $\{i,i+2\}$. A signing assigns a value in
$\{\pm1\}$ to each edge. Its signed adjacency matrix $A_\sigma$ is real
symmetric, and we write
\[
 \rho(A_\sigma)=\max\{|\lambda|:\lambda\in\spec(A_\sigma)\},
 \qquad \mu_n=\min_\sigma\rho(A_\sigma).
\]
Both spectral edges matter in this minimization. An upper bound for
$\lambda_{\max}$ alone does not give the required upper bound for $\rho$.

The earlier twisted-optimality conjecture \cite{Suvagiya2026old} compared
$\mu_n$ with the benchmark
\begin{equation}
 \tw(n)^2=4+2\cos\frac{2\pi}{n}+2\cos\frac{4\pi}{n}
 \qquad(n\text{ even},\ n\ge8).
 \label{eq:twisted}
\end{equation}
That preprint has been withdrawn and incorporated into
\cite{Suvagiya2026}. The revised Theorem 26 gives a period-eight signing
with radius
\begin{equation}
 \rstar=\sqrt{4+\sqrt{10+2\sqrt5}}
       =2.7936044933\ldots
 \label{eq:rstar}
\end{equation}
on $8m$ vertices for $m\ge4$, and Conjecture 28 asserts $\mu_{8m}=\rstar$
for $m\ge4$. The construction has positive Hamilton-cycle holonomy;
Remark 27 identifies the opposite holonomy as absent from the reported
search. Changing that holonomy while preserving the triangle signs gives
a strict improvement at every finite size.

Fix the word
\begin{equation}
 t=(1,1,-1,1,-1,-1,1,-1).
 \label{eq:word}
\end{equation}
For $n=8m$, begin with positive step-one edges and step-two signs
$t_{i\bmod8}$. Let $A_m^+$ be this adjacency matrix. Let $A_m^-$ be obtained
by reversing the signs of exactly three edges:
\begin{equation}
 \{n-1,0\},\qquad \{n-2,0\},\qquad \{n-1,1\}.
 \label{eq:seam}
\end{equation}
These are distinct edges for every $n\ge8$.

\begin{theorem}\label{thm:main}
For every integer $m\ge1$, define
\[
 R(s)=\sqrt{4+\sqrt{8+s+\sqrt{26-3s}}}\qquad(-2\le s\le2).
\]
Then
\begin{equation}
 \rho(A_m^+)=R(2)=\rstar,
 \qquad
 \rho(A_m^-)=R\!\left(2\cos\frac{\pi}{m}\right)<\rstar.
 \label{eq:main}
\end{equation}
Moreover, among all signings with the prescribed labeled triangle signs
$t_{i\bmod8}$, the negative-holonomy switching class is the unique
minimizing switching class. In particular,
\[
 \mu_{8m}\le R\!\left(2\cos\frac{\pi}{m}\right)<\rstar
 \qquad(m\ge4),
\]
so Conjecture 28 of \cite{Suvagiya2026} is false as stated.
\end{theorem}

The uniqueness assertion is confined to the prescribed triangle signs.
The theorem supplies no matching lower bound for arbitrary triangle words,
and therefore no classification of unrestricted minimizers. The dispersion relation and exact finite formula are recorded in the
earlier period-eight research manuscript \cite{EarlierPeriodEight}.
The present application concerns the revised conjecture and makes explicit
why its proposed endpoint is missed by every finite negative-holonomy grid.

A separate issue is to obtain a short proof for an explicit family outside
the $8$-divisible subsequence. The following statement is independent of
Theorem~\ref{thm:main}.

\begin{theorem}\label{thm:r2}
Let $k\ge6$, $n=8k+2$, and let $B_n$ be the signed adjacency matrix with
all step-one signs positive and the step-two sign word
\[
 \tau=t^k\mathbin{\|}(1,-1).
\]
Then
\begin{equation}
 \frac{198}{25}I_n-B_n^2\succ0,
 \qquad
 \rho(B_n)^2<\frac{198}{25}<\tw(n)^2.
 \label{eq:r2}
\end{equation}
\end{theorem}

Appendix~\ref{app:r2} gives the full analytic argument and its finite
rational premises. The latter have been checked both through the block
recurrence and independently from the signed graph by scalar Schur
elimination. The theorem is thus a finite-certificate-assisted analytic
result. The signing is inherited from the earlier residue-two construction;
the increment is the analytic closure and its corrected exact certificates.
The proof does not rely on the historical all-even classification.

The character interpretation of twists is part of established spectral
methods. Luo and Roy \cite{LuoRoy2026}, particularly their Theorems 3.6 and
3.13, study graph spectra over homology-character families and endpoint
attainment. Their setting provides context for holonomy, but does not give
an all-signing minimum for this fixed cycle-square support. The
parity-closed-walk identity of Chen, van Dam, and Bu
\cite[Theorem 3.1]{ChenDamBu2024} expresses signed moments averaged over
all signings; an average is likewise not a uniform lower bound for each
signing. Our proofs instead use the explicit fiber polynomial for one
triangle word and exact positive-definiteness certificates for another.

\section{Switching and finite holonomy}

Write $a_i$ for the sign of $\{i,i+1\}$ and $d_i$ for the sign of
$\{i,i+2\}$, with cyclic indices. The triangle sign and Hamilton-cycle
holonomy are
\[
 \tau_i=a_i a_{i+1}d_i,\qquad \alpha=\prod_{i=0}^{n-1}a_i.
\]
Switching by vertex signs $g_i\in\{\pm1\}$ replaces $A$ by $GAG$,
where $G=\diag(g_i)$. It preserves the spectrum, every $\tau_i$, and
$\alpha$.

\begin{lemma}\label{lem:gauge}
For prescribed triangle signs $(\tau_i)_{i\in\Z/n\Z}$, there are exactly
two switching classes, indexed by $\alpha\in\{\pm1\}$. Each has a unique
representative with
\[
 a_0=\cdots=a_{n-2}=1,\qquad a_{n-1}=\alpha,
 \qquad d_i=\tau_i a_i a_{i+1}.
\]
\end{lemma}
\begin{proof}
Set $g_0=1$ and recursively $g_{i+1}=g_i a_i$ for $0\le i<n-1$.
After switching, the first $n-1$ step-one signs are positive. Their product
with the remaining sign is $\alpha$, so that remaining sign is $\alpha$.
The triangle equations force every step-two sign. Conversely, the displayed
signs realize the prescribed data for either choice of $\alpha$.
A switch between two such representatives is constant along the path
$0,1,\ldots,n-1$, hence constant on all vertices and acts trivially.
Finally, distinct values of the switching invariant $\alpha$ cannot be
switching equivalent.
\end{proof}

For the repeated word \eqref{eq:word}, this gauge yields exactly $A_m^+$
and $A_m^-$ in \eqref{eq:seam}. Reversing only the Hamilton edge would
change two triangle signs. The two additional step-two reversals are
therefore essential.

Extend $u\in\C^{8m}$ to the integer line by
\begin{equation}
 u_{i+8m}=\alpha u_i.
 \label{eq:extension}
\end{equation}
Let $t_i=t_{i\bmod8}$. In either seam gauge the adjacency action is the
restriction of
\begin{equation}
 (Lu)_i=u_{i-1}+u_{i+1}+t_{i-2}u_{i-2}+t_i u_{i+2}.
 \label{eq:infinite-action}
\end{equation}
In particular, the signs of all cut-crossing edges in the finite matrix
agree with \eqref{eq:extension}, including for $m=1$.

\begin{lemma}\label{lem:bloch}
The complexification of $A_m^\alpha$ is unitarily equivalent to
\[
 \bigoplus_{z^m=\alpha}H(z),
\]
where, in the order $0,1,\ldots,7$,
\begin{equation}
 H(z)=\begin{pmatrix}
0&1&1&0&0&0&z^{-1}&z^{-1}\\
1&0&1&1&0&0&0&-z^{-1}\\
1&1&0&1&-1&0&0&0\\
0&1&1&0&1&1&0&0\\
0&0&-1&1&0&1&-1&0\\
0&0&0&1&1&0&1&-1\\
z&0&0&0&-1&1&0&1\\
z&-z&0&0&0&-1&1&0
\end{pmatrix}.
\label{eq:fiber}
\end{equation}
\end{lemma}
\begin{proof}
For each root $z^m=\alpha$ and $v\in\C^8$, set
$u_{8j+a}=m^{-1/2}z^jv_a$ for $0\le j<m$ and $0\le a<8$.
This satisfies \eqref{eq:extension}. Substitution into
\eqref{eq:infinite-action} gives \eqref{eq:fiber}. For two different
allowed phases $z,w$, the ratio $z\overline w$ is a nontrivial $m$th
root of unity, so
\[
 \sum_{j=0}^{m-1}(z\overline w)^j=0.
\]
Thus the $m$ fiber spaces are mutually orthogonal, each has dimension
eight, and together they exhaust dimension $8m$. The displayed map is
unitary. Since $|z|=1$, the fiber matrix is Hermitian.
\end{proof}

\section{The exact antiperiodic spectrum}

\begin{lemma}\label{lem:poly}
For $|z|=1$ and $s=z+z^{-1}$,
\begin{align}
 \det(xI_8-H(z))=P(x,s)
  &:=x^8-16x^6+80x^4-128x^2+38 \notag\\
  &\quad+s(-2x^4+16x^2-13)+s^2.
 \label{eq:poly}
\end{align}
Its eight roots, with all signs independently chosen, are
\begin{equation}
 \pm\sqrt{4\pm\sqrt{8+s\pm\sqrt{26-3s}}}.
 \label{eq:all-roots}
\end{equation}
They are all real, nonzero, and distinct for every $-2\le s\le2$.
\end{lemma}
\begin{proof}
Expanding the determinant of \eqref{eq:fiber} as a Laurent polynomial
in $z$ gives
\[
 x^8-16x^6+80x^4-128x^2+40
 +(z+z^{-1})(-2x^4+16x^2-13)+(z^2+z^{-2}).
\]
Using $z^2+z^{-2}=s^2-2$ proves \eqref{eq:poly}.
This is a finite polynomial identity and can also be checked by the exact
symbolic determinant in the accompanying verifier.

Put $y=x^2$. A translation gives
\begin{equation}
 P(x,s)=(y-4)^4-(16+2s)(y-4)^2+s^2+19s+38.
 \label{eq:centered}
\end{equation}
Solving this as a quadratic in $(y-4)^2$ gives
\[
 u_\pm(s)=8+s\pm\sqrt{26-3s}.
\]
The function $s^2+19s+38$ is increasing on $[-2,2]$ and has value $4$
at $-2$. Since $8+s>0$, the identity
\[
 (8+s)^2-(26-3s)=s^2+19s+38\ge4
\]
proves $u_->0$. Also $u_+>u_-$, and
\[
 u_+'(s)=1-\frac3{2\sqrt{26-3s}}>0,
 \qquad u_+(s)\le10+2\sqrt5<16.
\]
Consequently
\[
 0<4-\sqrt{u_+}<4-\sqrt{u_-}
   <4+\sqrt{u_-}<4+\sqrt{u_+}.
\]
These are the four distinct positive roots in $y$. Taking their positive
and negative square roots gives exactly \eqref{eq:all-roots}.
\end{proof}

\begin{proof}[Proof of Theorem~\ref{thm:main}]
Lemma~\ref{lem:poly} implies
\[
 \rho(H(z))=R(s),\qquad s=z+z^{-1}.
\]
The derivative computation in that lemma and strict monotonicity of the
outer square roots show that $R$ increases strictly on $[-2,2]$.
For positive holonomy, the grid $z^m=1$ contains $z=1$, so its maximal
value of $s$ is $2$. For negative holonomy, the grid is
\[
 z_j=\exp\!\left(\frac{(2j+1)\pi\mathrm i}{m}\right),
 \qquad 0\le j<m,
\]
and its maximal value of $s$ is $2\cos(\pi/m)$. Lemma~\ref{lem:bloch}
therefore proves both equalities in \eqref{eq:main}.
For every finite $m\ge1$, $\cos(\pi/m)<1$, which proves the strict
inequality. Lemma~\ref{lem:gauge} shows that these are exactly the two
switching classes for the prescribed triangle signs, so the strict
comparison proves the stated constrained uniqueness. Since $A_m^-$ is
an admissible real signing, the unrestricted upper bound follows.
\end{proof}

To identify \eqref{eq:rstar} with the algebraic number used in the revised
conjecture, put
\[
 f(x)=x^4-2x^3-6x^2+12x-4.
\]
An exact expansion gives $P(x,2)=f(x)f(-x)$. On $x\ge\sqrt6$,
\[
 f''(x)=12(x^2-x-1)>0,\qquad
 f'(\sqrt6)=12\sqrt6-24>0,\qquad f(\sqrt6)=-4.
\]
Hence $f$ has exactly one root on that ray. For $x>\sqrt6$,
$f(-x)-f(x)=4x(x^2-6)>0$. The largest positive root of $P(x,2)$ is
$R(2)>\sqrt6$. It must be a root of $f$, because $f(-R(2))=0$ would
force $f(R(2))<0$ and then a still larger positive root of $f$, contrary
to maximality. Thus $R(2)$ is the largest real root of $f$.

\begin{corollary}\label{cor:asymptotic}
The negative-holonomy radii increase strictly to $\rstar$ as $m$ increases,
and
\begin{equation}
 \rstar-\rho(A_m^-)
 =\frac{\pi^2}{4\rstar\sqrt{10+2\sqrt5}}
   \left(1-\frac3{4\sqrt5}\right)m^{-2}+O(m^{-4}).
 \label{eq:asymptotic}
\end{equation}
\end{corollary}
\begin{proof}
The sequence $2\cos(\pi/m)$ increases strictly to $2$, and $R$ is strictly
increasing and smooth near $2$. Since
\[
 2\cos(\pi/m)=2-\pi^2m^{-2}+O(m^{-4}),
 \qquad
 R'(2)=\frac{1-3/(4\sqrt5)}{4\rstar\sqrt{10+2\sqrt5}},
\]
Taylor's theorem proves \eqref{eq:asymptotic}.
\end{proof}

There is no parity exception: when $m$ is odd, the allowed phase $z=-1$
has $s=-2$ and cannot recover the omitted maximum at $s=2$.
At $m=1$ the single fiber is $H(-1)$, with
$\rho(A_1^-)=\sqrt{6+\sqrt2}$.

\section{An exact certificate at order 32}

At $m=4$, Theorem~\ref{thm:main} gives
\begin{equation}
 \rho(A_4^-)
 =\sqrt{4+\sqrt{8+\sqrt2+\sqrt{26-3\sqrt2}}}
 =2.7842697993\ldots.
 \label{eq:n32}
\end{equation}
The characteristic polynomial computed directly from the signed
$32\times32$ integer matrix is $Q(x)^2$, where
\begin{align*}
 Q(x)={}&x^{16}-32x^{14}+416x^{12}-2816x^{10}+10568x^8\\
       &-21632x^6+22168x^4-9408x^2+1262.
\end{align*}
This agrees with $P(x,\sqrt2)P(x,-\sqrt2)$; each value of $s$ occurs twice
on the antiperiodic four-cell grid.

A separate rational certificate avoids both the fiber calculation and
numerical eigenvalues. Every leading principal minor of each integer matrix
\[
 279I_{32}-100A_4^-,\qquad 279I_{32}+100A_4^-
\]
is strictly positive. The certificate records all $64$ integer minors;
fraction-free elimination checks every division exactly. Sylvester's
criterion yields $\rho(A_4^-)<279/100$. Meanwhile,
\[
 f(279/100)=-\frac{6766519}{100000000}<0,
 \qquad f(14/5)=\frac{76}{625}>0.
\]
The monotonicity proved above gives the exact separation
\[
 \rho(A_4^-)<2.79<\rstar<2.8.
\]
The direct graph certificate and the all-$m$ analytic proof provide
independent ways of checking the failure at $n=32$.

The next extremal question remains to determine $\mu_{8m}$ and its
minimizing switching classes. The period-eight negative-holonomy family is
an explicit upper bound, not a proof of optimality. Establishing optimality
would require a lower bound over arbitrary triangle words, including
nonperiodic words. Similarly, Theorem~\ref{thm:r2} is an existence theorem
for one signing on each order $8k+2$. Combining such witnesses does not
certify a complete truth set for the twisted benchmark without independent
lower bounds at the complementary orders.

\appendix
\section{A uniform residue-two Schur bound}\label{app:r2}

We prove Theorem~\ref{thm:r2}. In this appendix $n=8k+2=4\ell+2$,
where $\ell=2k$ is even and $k\ge6$. The letter $\ell$ counts four-vertex
blocks and is distinct from the cell count $m$ in the main text. Put
\[
 M_n=\frac{198}{25}I_n-B_n^2.
\]
All matrix inequalities in this appendix concern real symmetric matrices;
$X\succ0$ means positive definite. The norm $\norm{\cdot}_2$ is the
Euclidean operator norm, and $\norm{\cdot}_F$ is the Frobenius norm.

\subsection{The exact block recurrence}

Partition the vertices into
\[
 V_0=\{0,1\},\qquad
 V_j=\{2+4(j-1),\ldots,5+4(j-1)\}\quad(1\le j\le\ell).
\]
The diagonal block on each $V_j$, $j\ge1$, is
\[
 D=\begin{pmatrix}
 98/25&0&-1&0\\0&98/25&0&-1\\-1&0&98/25&0\\0&-1&0&98/25
 \end{pmatrix}.
\]
The block $M_n[V_j,V_{j+1}]$ equals $E_+$ for odd $j$ and $E_-$ for
even $j$, where
\[
 E_+=\begin{pmatrix}-1&0&0&0\\0&1&0&0\\-1&2&1&0\\2&-1&0&-1\end{pmatrix},
 \qquad
 E_-=\begin{pmatrix}-1&0&0&0\\0&1&0&0\\-1&-2&1&0\\-2&-1&0&-1\end{pmatrix}.
\]
The exceptional blocks are
\begin{gather*}
 G_0=M_n[V_0,V_0]=(98/25)I_2,\qquad H_0=D,\\
 R_0=M_n[V_0,V_1]=\begin{pmatrix}-1&-2&1&0\\-2&-1&0&-1\end{pmatrix},\\
 C_0=M_n[V_0,V_\ell]=\begin{pmatrix}-1&0&-1&0\\0&-1&0&-1\end{pmatrix},
 \qquad
 W_0=M_n[V_1,V_\ell]=\begin{pmatrix}0&0&-1&0\\0&0&0&1\\0&0&0&0\\0&0&0&0\end{pmatrix}.
\end{gather*}
All other nonadjacent bulk blocks vanish. To verify these identities at
all lengths, the only nonzero entries of $B_n^2$ on the forward cyclic
diagonals, apart from transposes, are
\begin{align*}
 (B_n^2)_{i,i}&=4,&
 (B_n^2)_{i,i+1}&=\tau_{i-1}+\tau_i,&
 (B_n^2)_{i,i+2}&=1,\\
 (B_n^2)_{i,i+3}&=\tau_i+\tau_{i+1},&
 (B_n^2)_{i,i+4}&=\tau_i\tau_{i+2}.&&
\end{align*}
They follow by enumerating the possible two-edge walks. Since $n\ge50$,
these cyclic offsets do not collide. Substitution of $t^k\|(1,-1)$ gives
the displayed blocks without a length-dependent assumption.

Let $E_j=E_+$ for even $j$ and $E_j=E_-$ for odd $j$. Set $X_0=D$.
Before closing the right boundary, elimination of the current four-vertex
pivot gives the length-independent recurrence
\begin{align}
 X_{j+1}&=D-E_j^TX_j^{-1}E_j,&
 R_{j+1}&=-R_jX_j^{-1}E_j,\notag\\
 W_{j+1}&=-E_j^TX_j^{-1}W_j,&
 G_{j+1}&=G_j-R_jX_j^{-1}R_j^T,\label{eq:recurrence}\\
 H_{j+1}&=H_j-W_j^TX_j^{-1}W_j,&
 C_{j+1}&=C_j-R_jX_j^{-1}W_j.\notag
\end{align}
The terminal elimination index is $p=\ell-2$, which is even.
At that step only, the coupling to the retained last block becomes
$W_p+E_+$. Thus the normalized six-dimensional core is
\begin{equation}
 S_\ell=
 \begin{pmatrix}
 G_p-R_pX_p^{-1}R_p^T&C_p-R_pX_p^{-1}(W_p+E_+)\\
 *&H_p-(W_p+E_+)^TX_p^{-1}(W_p+E_+)
 \end{pmatrix}.
 \label{eq:core}
\end{equation}
Here $*$ denotes the transpose block. Repeated Schur congruence proves
that $M_n\succ0$ if all eliminated pivots and $S_\ell$ are positive.
The subscript of $S_\ell$ is a block count: in particular, $S_{26}$
corresponds to graph order $106$.

\subsection{Finite exact premises}

Define $F_\pm(X)=D-E_\pm^TX^{-1}E_\pm$ and
$\Phi=F_-\circ F_+$. Put $Z=\Phi^{12}(D)=X_{24}$,
$Y=F_+(Z)$, and $L_Z=Z^{-1}E_+Y^{-1}E_-$. Use the rational matrices
\begin{align*}
 P&=\frac1{10000}\begin{pmatrix}
10766&87&19&974\\87&12664&148&-2418\\19&148&10093&-25\\974&-2418&-25&14009
\end{pmatrix},\\
 Q&=\frac1{10000}\begin{pmatrix}
11503&614&990&-1101\\614&10470&15&113\\990&15&12299&-2632\\-1101&113&-2632&13260
\end{pmatrix}.
\end{align*}
Set $r=10^{-10}$. The following finite inequalities are the complete
certificate premises used below:
\begin{enumerate}
\item $X_0,\ldots,X_{24}\succ0$ and $Z,Y\succ I_4/2$;
\item $(9/10)I_4\prec P,Q\prec2I_4$;
\item $L_Z^TPL_Z\prec(2/5)P$ and $L_ZQL_Z^T\prec(2/5)Q$;
\item $\norm{\Phi(Z)-Z}_F<r/40$;
\item $R_{24}QR_{24}^T\prec10^{-10}I_2$ and
$W_{24}^TQW_{24}\prec10^{-10}I_4$;
\item $S_{26}\succ(1/50)I_6$;
\item $S_\ell\succ0$ for $\ell=12,14,16,18,20,22,24$.
\end{enumerate}
They have been verified using exact rational arithmetic. For clarity, their
verification requires no limiting or approximate matrix: all matrices are
specified by the displayed rational initial data and a finite number of
recurrence steps. Positive definiteness is checked by successive scalar
Schur pivots. For a symmetric matrix $K$, start with $K^{(0)}=K$; at each
step require $d_j=K^{(j)}_{11}>0$ and replace the remaining block by
$K^{(j+1)}=K^{(j)}_{22}-K^{(j)}_{21}(d_j)^{-1}K^{(j)}_{12}$.
All resulting scalars are rational, so each test is exact. The residual
test in item 4 is the rational inequality
\[
 \sum_{a,b}(\Phi(Z)-Z)_{ab}^2<(r/40)^2.
\]
The accompanying certificate records the matrices and the positive pivots.
An independent replay constructs $B_{106}$ from its edges, eliminates
$96$ scalar vertices to recover the entrance data, and verifies items
1--5 without generating that entrance by the block recurrence. Direct
scalar elimination from the graphs also verifies the seven base orders
$50,58,66,74,82,90,98$ and the normalized seed at $106$.

\subsection{A local Riccati contraction}

For symmetric $U$, define the order-unit norm
\[
 |U|_P=\inf\{a\ge0:-aP\preceq U\preceq aP\}
      =\norm{P^{-1/2}UP^{-1/2}}_2.
\]
Consider the closed ball
$\mathcal K=\{X=X^T:|X-Z|_P\le r\}$. Since $P\prec2I$,
\[
 \norm{X-Z}_2\le2r=:e.
\]
The certified lower bound $Z\succ I/2$ gives $X\succ I/3$ and
$\norm{X^{-1}}_2\le3$. The inverse identity yields
\[
 \norm{X^{-1}-Z^{-1}}_2\le6e.
\]
Both $E_\pm$ have squared Frobenius norm $14$. Writing $Y_X=F_+(X)$,
we obtain
\[
 \norm{Y_X-Y}_2\le84e<1/6,
 \qquad Y_X\succ I/3,
 \qquad \norm{Y_X^{-1}-Y^{-1}}_2\le504e.
\]
Define $L(X)=X^{-1}E_+Y_X^{-1}E_-$. Splitting its difference from
$L_Z$ into two terms gives
\begin{equation}
 \norm{L(X)-L_Z}_2
 \le14(6\cdot3+2\cdot504)e=14364e.
 \label{eq:Lperturb}
\end{equation}
For the two transfer norms
\[
 \norm{L}_{P,\mathrm{col}}=\norm{P^{1/2}LP^{-1/2}}_2,
 \qquad
 \norm{L}_{Q,\mathrm{row}}=\norm{Q^{-1/2}LQ^{1/2}}_2,
\]
conversion from the Euclidean norm costs at most
$\sqrt{20/9}<3/2$. Therefore the perturbation in either norm is less
than $21546e=43092r<10^{-4}$. The certified bounds at $Z$, together with
$\sqrt{2/5}+10^{-4}<2/3$, imply throughout $\mathcal K$ that
\begin{equation}
 L(X)^TPL(X)\prec\frac49P,
 \qquad L(X)QL(X)^T\prec\frac49Q.
 \label{eq:dual-contractions}
\end{equation}
Direct differentiation gives
$D\Phi_X[U]=L(X)^TUL(X)$. If $-aP\preceq U\preceq aP$, the first
inequality in \eqref{eq:dual-contractions} bounds this derivative in
$|\cdot|_P$ by $4a/9$. Integrating along a line segment in the convex
ball proves that $\Phi$ is $4/9$-Lipschitz there.

The residual premise gives
$|\Phi(Z)-Z|_P<(10/9)(r/40)=r/36$. Consequently
\[
 |\Phi(X)-Z|_P<r/36+(4/9)r=(17/36)r<r.
\]
The ball is invariant. Banach's fixed-point theorem supplies a unique
$X_*\in\mathcal K$ with $\Phi(X_*)=X_*$. The actual even trajectory
starts at $Z$, so, with $\theta=4/9$,
\begin{equation}
 \norm{X_{24+2h}-X_*}_2\le4r\theta^h
 \qquad(h\ge0).
 \label{eq:pivot-decay}
\end{equation}
The constant $4r$ follows already from the diameter of the ball. All these
even pivots and their intermediate odd partners exceed $I/3$. Together
with the finite entrance this proves positivity of every pivot needed
at any length.

\subsection{Response decay}

Two consecutive updates in \eqref{eq:recurrence} act by
$R\mapsto RL(X)$ and $W\mapsto L(X)^TW$. For these two orientations,
respectively, use
\[
 \alpha(R)=\norm{RQ^{1/2}}_2,
 \qquad \beta(W)=\norm{Q^{1/2}W}_2.
\]
The second inequality in \eqref{eq:dual-contractions} contracts both
quantities by $q=2/3$. Hence the entrance premises imply
\begin{align*}
 R_{24+2h}QR_{24+2h}^T&\prec10^{-10}q^{2h}I_2,\\
 W_{24+2h}^TQW_{24+2h}&\prec10^{-10}q^{2h}I_4.
\end{align*}
Since $Q\succ(9/10)I$, it follows that
\begin{equation}
 \norm{R_{24+2h}}_2,\ \norm{W_{24+2h}}_2<a q^h,
 \qquad a=1/30000.
 \label{eq:even-response}
\end{equation}
Each one-step transfer has norm at most $3\sqrt{14}<12$. Thus both
members of the $h$th pair obey
\begin{equation}
 \norm{R_j}_2,\ \norm{W_j}_2\le bq^h,
 \qquad b=12a=1/2500,
 \quad j\in\{24+2h,25+2h\}.
 \label{eq:pair-response}
\end{equation}
The row orientation $LQL^T$ is needed here. A bound on $L^TQL$ alone
would not justify these response estimates.

\subsection{The limiting core and its uniform tail}

Use the actual infinite trajectory, not a trajectory with constant pivots,
to define the absolutely convergent series
\begin{align*}
 G_\infty&=G_0-\sum_{j\ge0}R_jX_j^{-1}R_j^T,\\
 H_\infty&=H_0-\sum_{j\ge0}W_j^TX_j^{-1}W_j,\\
 C_\infty&=C_0-\sum_{j\ge0}R_jX_j^{-1}W_j.
\end{align*}
Convergence follows from \eqref{eq:pair-response} and
$\norm{X_j^{-1}}_2\le3$ after $j=24$. Set
\[
 S_\infty=\begin{pmatrix}
 G_\infty&C_\infty\\
 C_\infty^T&H_\infty-E_+^TX_*^{-1}E_+
 \end{pmatrix}.
\]
Take even $\ell\ge26$ and write its terminal index as
$p=\ell-2=24+2h$. Each of the three Schur-series suffixes after index
$p$ is bounded in norm by
\begin{equation}
 T_h=\frac{6b^2q^{2h}}{1-q^2}.
 \label{eq:Th}
\end{equation}
Indeed, every increment in the $j$th pair has norm at most
$3b^2q^{2j}$, and there are two increments per pair. The bound in
\eqref{eq:Th} harmlessly includes the term at $p$ as well.

The terminal linear terms from \eqref{eq:core} satisfy
\begin{align*}
 \norm{R_pX_p^{-1}E_+}_2&\le12a q^h,\\
 \norm{W_p^TX_p^{-1}E_++E_+^TX_p^{-1}W_p}_2&\le24a q^h.
\end{align*}
The inverse identity and \eqref{eq:pivot-decay}, with $\delta=4r$,
give
\[
 \norm{X_p^{-1}-X_*^{-1}}_2\le9\delta\theta^h,
 \qquad
 \norm{E_+^T(X_p^{-1}-X_*^{-1})E_+}_2
 \le144\delta\theta^h.
\]
For the differences of the three blocks in $S_\ell-S_\infty$ we
therefore have
\begin{align*}
 \norm{\Delta G}_2&\le T_h,\\
 \norm{\Delta C}_2&\le T_h+12a q^h,\\
 \norm{\Delta H}_2&\le T_h+24a q^h+144\delta\theta^h.
\end{align*}
The block estimate
\[
 \norm{\begin{pmatrix}U&V\\V^T&Z\end{pmatrix}}_2
 \le\norm U_2+2\norm V_2+\norm Z_2
\]
now gives the explicit uniform tail bound
\begin{align}
 \norm{S_\ell-S_\infty}_2
 &\le4T_h+48a q^h+144\delta\theta^h\notag\\
 &\le\frac{251089}{156250000}<\frac1{500}.
 \label{eq:tail}
\end{align}
The rational middle expression is the value of this majorant at $h=0$;
each geometric factor decreases for $h\ge0$. Because the series were
defined along the actual trajectory, the common initial terms cancel
exactly. Only the terminal pure-pivot term requires the inverse-limit
comparison above.

\subsection{Completion of the spectral estimate}

The exact seed $S_{26}\succ I_6/50$ and two applications of
\eqref{eq:tail} imply, for every even $\ell\ge26$,
\[
 \norm{S_\ell-S_{26}}_2<2/500=1/250,
 \qquad S_\ell\succ\left(\frac1{50}-\frac1{250}\right)I_6
                  =\frac2{125}I_6.
\]
Both tail errors are required: the seed is a finite core, not the limiting
core. The remaining even block counts $12\le\ell<26$ are exactly the
seven certified bases. Thus every core and every eliminated pivot is
positive. Schur congruence proves $M_n\succ0$, and the spectral theorem
then gives $\rho(B_n)^2<198/25$.

This use of Schur congruence asserts positivity only. The lower bound
$2/125$ for the reduced core is not a lower bound for the least Euclidean
eigenvalue of the original matrix $M_n$.
Finally, $\cos x>1-x^2/2$ for $x>0$ and $\pi^2<10$ yield
\[
 \tw(n)^2>8-\frac{20\pi^2}{n^2}
          >8-\frac{200}{n^2}\ge\frac{198}{25}
 \qquad(n\ge50).
\]
This completes the proof of Theorem~\ref{thm:r2}.

\paragraph{Reproducibility.}
\begingroup\raggedright
The accompanying source package includes
\nolinkurl{certificates/analytic/verify_r2_certificate.py} and
\nolinkurl{r2_exact_certificate.json}, together with the independent
\nolinkurl{replay_direct_graph_seed.py} and
\nolinkurl{replay_local_from_direct_graph.py} under
\nolinkurl{certificates/audit/}. For the separate antiperiodic check, the script
\nolinkurl{verify_antiperiodic_counterexample.py} under
\nolinkurl{certificates/new_conjecture/} constructs the graph,
checks the symbolic identities, and verifies all recorded integer minors.
The package README lists the exact execution commands and dependencies.
\par\endgroup


\begin{thebibliography}{9}
\bibitem{Suvagiya2026old}
V. Suvagiya,
\emph{Signed circulants at the Ramanujan bound},
arXiv:2607.18334v1, 19 July 2026;
withdrawn in version 2, 22 September 2026, and merged into
\cite{Suvagiya2026}.
\url{https://arxiv.org/abs/2607.18334}.

\bibitem{Suvagiya2026}
V. Suvagiya,
\emph{Parity families and signed spectra: kernel averaging,
near-Ramanujan bounds, and exact circulant models},
arXiv:2607.17343v2, 22 September 2026,
Theorem 26, Remark 27, and Conjecture 28.
\url{https://arxiv.org/abs/2607.17343v2}.

\bibitem{LuoRoy2026}
Y. Luo and A. Roy,
\emph{Homological spectral graph theory and weighted cycle counting},
arXiv:2403.01550v4, 4 August 2026.
\url{https://arxiv.org/abs/2403.01550v4}.

\bibitem{ChenDamBu2024}
L. Chen, E. R. van Dam, and C. Bu,
Spectra of power hypergraphs and signed graphs via parity-closed walks,
\emph{Journal of Combinatorial Theory, Series A} \textbf{207} (2024),
105909.
\url{https://doi.org/10.1016/j.jcta.2024.105909}.

\bibitem{EarlierPeriodEight}
\emph{The exact period-eight phase}, Section 4 of the earlier period-eight
research manuscript, frozen repository source at commit
\texttt{085ea698475b7b32e0ae57457ec903a922248f69}.
\href{https://github.com/whzy3185/math/blob/085ea698475b7b32e0ae57457ec903a922248f69/research/paper_strengthening/manuscript_period8_jgt/sections_en/04_period8_exact.tex}{Archived source section}.

\end{thebibliography}
\end{document}
